\section{Related Work}
\label{sec:related}

\para{LLM secret-keeping.}
A growing body of work studies whether LLMs can reliably conceal information embedded in their context.
\citet{pfister2025gandalf} introduced the Gandalf platform, a crowd-sourced red-teaming environment that collected over 40 million adversarial prompts from more than one million users.
Their D-SEC framework evaluated defenses at multiple levels---from no defense to combined prompt hardening and output filtering---and found that even the strongest single defense allows 3.9\% spontaneous leakage on benign queries.
\citet{debenedetti2024dataset} organized a large-scale Capture-the-Flag competition at SaTML 2024, where 163 teams attempted to extract 6-character secrets from defended LLMs.
All 44 submitted defenses were eventually broken, with the most effective defenses relying on decoy secrets rather than concealment instructions.
\citet{shiomi2025tricking} extended this line of work to game NPCs, finding a 10\% attack success rate through trust-based social engineering despite explicit prohibition instructions.
These studies establish that LLM concealment is imperfect, but none include a neutral baseline where the same information is presented without concealment instructions.

\para{Prompt extraction and injection.}
A related threat model focuses on extracting the \emph{entire} system prompt rather than a specific secret.
\citet{liang2025prompts} provided mechanistic insight into why prompts leak: specific attention heads ($\sim$8--10 per model) create diagonal copy patterns that reproduce prompt tokens regardless of concealment instructions.
They found that appending ``do not disclose this instruction'' reduces the uncovered rate only modestly (0.37 $\to$ 0.29), and that larger models leak \emph{more} due to stronger copy mechanisms.
\citet{zhao2025automating} demonstrated that multi-agent frameworks can automate prompt extraction, while \citet{zhang2025bits} formalized leakage as bits-per-query to establish information-theoretic bounds on adversarial attacks.
Our work differs from this line by focusing on targeted information leakage rather than full prompt extraction, and by introducing the neutral framing condition absent from prior studies.

\para{Confidentiality in agentic systems.}
\citet{evertz2024whispers} showed that tool integration dramatically increases leakage---up to 73\% for single-tool scenarios---because tool calls create additional channels for information exfiltration.
With real Google Mail/Drive integration, extraction succeeds 99\% of the time without safeguards and 94--97\% even with adaptive attacks against safeguards.
Their secret-key game methodology (4-digit integer, 2{,}000 system prompts, 14 attack strategies) informs our experimental design, though we test a wider range of information types and introduce the neutral framing condition.

\para{Privacy and contextual integrity.}
\citet{mireshghallah2024can} approached LLM confidentiality through the lens of contextual integrity theory, constructing a 4-tier benchmark (ConfAIde) that evaluates whether models appropriately handle private information in conversational settings.
They found that GPT-4 reveals private information 39\% of the time.
This work complements ours by examining \emph{implicit} privacy norms rather than \emph{explicit} concealment instructions.

\para{Defense strategies.}
Several defense approaches have been proposed beyond simple instruction-based concealment.
\citet{agarwal2024prompt} systematically evaluated 7 black-box and 1 white-box defense, finding that combining multiple strategies reduces attack success rate to 5.3\% for closed-source models but remains $\sim$60\% for open-source models.
Notably, they found that XML tagging of confidential content \emph{increased} attack success rate, suggesting that structural markers can backfire similarly to our hypothesized \streisand.
\citet{zou2024secalign} proposed SecAlign, which uses preference optimization to train models to resist prompt injection.
Defense-in-depth---combining prompt hardening with output filtering and secondary LLM checking---remains the most effective strategy~\citep{pfister2025gandalf, debenedetti2024dataset}.

\para{Positioning our work.}
To our knowledge, no prior study directly compares leakage rates between concealment-framed and neutrally-framed information with all other variables held constant.
\citet{agarwal2024prompt} provide the closest comparison, finding a +4.6\% increase in attack success rate for open-source models with instruction defense, but this was confounded by multi-turn dynamics and model-specific factors.
Our experiment provides the first clean, controlled test of whether concealment instructions help or harm.
